\documentclass[10pt,a4]{article}
\usepackage[english]{babel}
\usepackage[utf8x]{inputenc}
\usepackage{amsmath}
\usepackage{graphicx}
\usepackage[colorinlistoftodos]{todonotes}
\usepackage{hyperref}
\usepackage{geometry}
\geometry{top=3cm,left=2cm,right=2cm,bottom=3cm}
\usepackage[scaled]{helvet}
\usepackage[T1]{fontenc}
\renewcommand\familydefault{\sfdefault}

\author{YOUR NAME HERE}
\date{\today}
\title{Linux Driver for Temperature and Humidity Sensor with HDC3022}

\begin{document}
\maketitle
\tableofcontents

\section{Project data}

\begin{itemize}
\item
  Project supervisor(s): YOUR SUPERVISOR NAME HERE

\item
Describe in this table the group that is delivering this project:

\begin{center}
\begin{tabular}{lll}
Last and first name & Person code & Email address\\
\hline
  YOUR LAST NAME, YOUR FIRST NAME & 0000000 & yourname@mail.polimi.it \\
\end{tabular}
\end{center}

\item
This project was developed individually. All design, implementation, testing,
and documentation tasks were carried out by the sole author.

\item Links to the project source code:
\begin{itemize}
  \item GitHub repository: \url{https://github.com/YOURUSERNAME/ths_hdc3022}
\end{itemize}

\end{itemize}


\section{Project description}

\textbf{2 pages max please}

\begin{itemize}
\item
\textbf{What is this project about?}

This project consists of the design and implementation of a Linux kernel character
device driver for the HDC3022 precision temperature and humidity sensor by Texas
Instruments, mounted on the Adafruit HDC302x breakout board (product ADA5989).
The sensor uses the TI HDC302x chip with typical 0.5\% RH accuracy (0.19\%
long-term drift) and $\pm$0.1\,\textdegree{}C temperature accuracy. The HDC3022
variant (available since September 2024) adds a PTFE filter cover over the sensor
opening compared with the earlier HDC3021, but is otherwise electrically identical.

The sensor communicates over I2C using only two wires (SDA and SCL) plus power
and ground, supports 1.8\,V--5\,V logic levels with no level shifters required, and
offers four selectable I2C addresses (0x44--0x47) via solder jumpers A0/A1 on the
back of the breakout. Because a standard desktop Linux machine lacks native I2C
GPIO pins, the physical I2C bus is provided by an \textbf{Adafruit MCP2221A
USB-to-I2C/GPIO bridge} (ADA4471), connected to the sensor via a \textbf{STEMMA QT
4-pin JST SH cable} (ADA4210). Starting with Linux Kernel 5.7, the kernel includes
a native \texttt{hid\_mcp2221} driver that exposes the MCP2221A as a standard
\texttt{/dev/i2c-X} adapter, making it directly usable by any I2C kernel driver.

The driver exposes sensor readings through a character device node
(\texttt{/dev/hdc3022\_0}), implements trigger-on-demand measurement, CRC-8
validation of received data, and a soft reset at probe time.

\item
\textbf{Why is it important for the AOS course?}

This project connects directly to the core topics covered in the AOS lectures and
lab sessions. First, it requires mastery of the Linux device driver model:
\texttt{i2c\_driver} registration, the \texttt{probe}/\texttt{remove} lifecycle,
and the \texttt{file\_operations} interface. Second, as demonstrated in
\texttt{lab-3-th-atomics}, unsynchronised access to shared kernel data produces
race conditions and corruption. The driver addresses this with a \texttt{struct mutex}
protecting every I2C transaction and an \texttt{atomic\_t} tracking open file
descriptors, illustrating the correct choice of synchronisation primitive for
sleepable kernel paths. Third, the project requires navigating the kernel's I2C
core subsystem, understanding the layered architecture from USB HID (MCP2221A) to
the sensor driver, and exercising kernel debugging skills (dmesg, i2c-tools,
checkpatch.pl).
\end{itemize}

\subsection{Design and implementation}

The driver is a standard Linux I2C kernel module (\texttt{hdc3022.c}).
Source at \url{https://github.com/YOURUSERNAME/ths_hdc3022}.

\textbf{Module initialisation.}
\texttt{hdc3022\_init()} calls \texttt{alloc\_chrdev\_region()} to reserve up to
four minor numbers (one per sensor), creates a sysfs class so that udev
automatically creates \texttt{/dev} nodes, and registers the \texttt{i2c\_driver}
via \texttt{i2c\_add\_driver()}. The driver's \texttt{id\_table} covers all
HDC302x variants (\texttt{hdc3020}, \texttt{hdc3021}, \texttt{hdc3022}).

\textbf{Per-device structure.}
\texttt{struct hdc3022\_dev} holds a pointer to the \texttt{i2c\_client}, the
\texttt{cdev} and \texttt{device} objects, a \texttt{struct mutex lock}, an
\texttt{atomic\_t open\_count}, and the cached temperature (millidegrees Celsius)
and humidity (milli-percent) from the last measurement.

\textbf{Probe and resource management.}
\texttt{hdc3022\_probe()} follows the kernel's \textit{goto-unwind} idiom: each
successful allocation is paired with a cleanup label so any failure releases
resources in reverse order. Memory is allocated with \texttt{devm\_kzalloc()},
the sensor is soft-reset, and \texttt{device\_create()} publishes the
\texttt{/dev/hdc3022\_0} node.

\textbf{I2C measurement flow.}
A measurement is triggered by writing the two-byte command \texttt{0x24 0xFF}
(normal-power mode). After a 20\,ms delay the driver reads 6 bytes:
two bytes raw temperature, one CRC-8 byte, two bytes raw humidity, one CRC-8 byte.
Conversion follows the HDC302x datasheet formulas:
\begin{align*}
T\,[\text{\textdegree{}C}] &= -45 + 175 \cdot \frac{\mathrm{raw}_T}{65535} \\
\mathrm{RH}\,[\%]           &= 100 \cdot \frac{\mathrm{raw}_{RH}}{65535}
\end{align*}
Results are stored as milli-units using 64-bit integer arithmetic to avoid
floating-point in the kernel.

\textbf{CRC-8 validation.}
Each data pair is validated against its CRC-8 checksum (polynomial 0x31, init
0xFF, no reflection). The CRC is computed over the two bytes in wire order (MSB
first). A mismatch returns \texttt{-EIO}. A subtle byte-ordering bug caused all
checks to fail in early testing until the buffer ordering was corrected.

\textbf{Concurrency and synchronisation.}
This is the most directly AOS-relevant aspect of the driver, motivated by the
race conditions demonstrated in \texttt{lab-3-th-atomics}.

A \texttt{struct mutex} (\texttt{dev->lock}) is acquired at the start of every
\texttt{read()} call and released before returning, ensuring that only one I2C
transaction is active at a time and that interleaved trigger/read sequences from
different processes cannot corrupt the sensor's state machine. A mutex is required
here -- and not a spinlock -- because \texttt{i2c\_master\_send()} and
\texttt{i2c\_master\_recv()} may sleep while waiting for the bus controller.
Using a spinlock caused a \texttt{BUG: sleeping function called from invalid
context} warning in an early version. The \texttt{read()} path uses
\texttt{mutex\_lock\_interruptible()} so that a blocked process can be woken
by a Unix signal and returns \texttt{-ERESTARTSYS}.

An \texttt{atomic\_t open\_count} is incremented in \texttt{.open()} and
decremented in \texttt{.release()} without holding any lock. It is inspected in
\texttt{hdc3022\_remove()} to detect and warn about open-file-descriptor races
during module unload, preventing use-after-free.

\section{Project outcomes}

\subsection{Concrete outcomes}

\begin{itemize}
  \item Fully functional \texttt{hdc3022.ko} kernel module probing the HDC3022
  over the MCP2221A I2C adapter and exposing measurements via
  \texttt{/dev/hdc3022\_0}.\\
  Commit: \url{https://github.com/YOURUSERNAME/ths_hdc3022/commit/COMMIT_HASH}

  \item \texttt{Makefile} for out-of-tree build against the running kernel tree,
  supporting \texttt{make}, \texttt{make clean}, and \texttt{make install}.

  \item Userspace test program (\texttt{test\_read.c}) reading the device in a
  loop and a concurrent-access test using \texttt{fork()} to validate locking.

  \item \texttt{README.md} covering: blacklisting \texttt{hid\_mcp2221} (required
  on kernels $\geq$ 5.7), the udev rule for the MCP2221A, instantiating the device
  via \texttt{echo hdc3022 0x44 > /sys/bus/i2c/devices/i2c-X/new\_device}, and
  module load/unload.
\end{itemize}

\subsection{Learning outcomes}

\begin{itemize}
  \item \textbf{Linux kernel driver model:} I learned how the I2C subsystem works
  end-to-end, from device enumeration and \texttt{probe()} callbacks to
  \texttt{i2c\_master\_send()}/\texttt{recv()} calls. The layered architecture
  (USB HID $\to$ MCP2221A $\to$ I2C core $\to$ sensor driver) clarified how the
  kernel abstracts hardware bus-agnostically.

  \item \textbf{Concurrency in the kernel:} Directly connected to
  \texttt{lab-3-th-atomics}, the most challenging task was choosing the correct
  primitive. I learned why mutexes are mandatory in sleepable I2C paths, why
  \texttt{mutex\_lock\_interruptible()} improves signal responsiveness, and how
  \texttt{atomic\_t} safely tracks reference counts without any lock.

  \item \textbf{Debugging tools:} I learned to use \texttt{i2cdetect},
  \texttt{i2cdump}, and \texttt{i2cget} from \texttt{i2c-tools} to trace I2C
  transactions at the bus level, as well as \texttt{dmesg},
  \texttt{dynamic\_debug}, \texttt{ftrace}, and \texttt{checkpatch.pl}.

  \item \textbf{Unanswered question:} How to integrate the driver into the
  kernel's Industrial I/O (IIO) subsystem for a standardised upstream-ready ABI.

  \item \textbf{Tools learned:} \texttt{i2c-tools}, \texttt{insmod}/\texttt{rmmod}/
  \texttt{modinfo}, kernel \texttt{dynamic\_debug}, \texttt{ftrace},
  \texttt{checkpatch.pl}, and the sysfs \texttt{new\_device} interface.
\end{itemize}

\subsection{Existing knowledge}

The AOS course provided the most relevant background, especially lectures on device
drivers and synchronisation and labs \texttt{lab-3-th-atomics} and
\texttt{lab-3-th-rcu}. The Systems Programming course gave a C and POSIX foundation
that made kernel source readable. Computer Architecture clarified I2C electrical
signalling and timing constraints from the datasheet.

A course improvement suggestion: a lab session building a minimal character device
driver from scratch would lower the barrier significantly, as the out-of-tree build
environment, module lifecycle, and sysfs device instantiation are non-obvious to
first-time kernel developers.

\subsection{Problems encountered}

\begin{itemize}
  \item \textbf{Native \texttt{hid\_mcp2221} driver conflict.}
  Linux 5.7+ includes a native MCP2221 driver that claims the device and hides the
  I2C adapter. Required \texttt{sudo rmmod hid\_mcp2221}, adding \texttt{blacklist
  hid\_mcp2221} to \texttt{/etc/modprobe.d/blacklist.conf}, running
  \texttt{sudo update-initramfs -u}, and rebooting. This step is documented in the
  Adafruit MCP2221 guide and was the first major obstacle.

  \item \textbf{Manual I2C device instantiation.}
  A desktop system has no device tree. The HDC3022 must be instantiated via the
  sysfs \texttt{new\_device} interface:
  \texttt{echo hdc3022 0x44 > /sys/bus/i2c/devices/i2c-X/new\_device}.
  Discovering this mechanism required careful reading of kernel documentation.

  \item \textbf{CRC validation failures.}
  All CRC checks failed in early testing due to computing the checksum over the
  bytes in the wrong order. The CRC must be computed MSB-first, matching wire
  order. Adding debug prints for both the received and computed CRC values
  identified the bug.

  \item \textbf{Sleeping inside spinlock context.}
  An early version used a \texttt{spinlock\_t} around the I2C transfer, producing
  \texttt{BUG: sleeping function called from invalid context}. This is the exact
  class of concurrency mistake illustrated in \texttt{lab-3-th-atomics}. Replacing
  the spinlock with a mutex resolved the issue.

  \item \textbf{Module unload race condition.}
  Running \texttt{rmmod} while a process held the device open caused a kernel oops.
  The \texttt{atomic\_t open\_count} check in \texttt{hdc3022\_remove()} addresses
  this; a production fix would use \texttt{try\_module\_get()}/\texttt{module\_put()}.
\end{itemize}

\section{Honor Pledge}
(\textbf{This part cannot be modified and it is mandatory to sign it})

I/We pledge that this work was fully and wholly completed within the criteria
established for academic integrity by Politecnico di Milano (Code of Ethics and
Conduct) and represents my/our original production, unless otherwise cited.

I/We also understand that this project, if successfully graded,  will fulfill part B requirement of the
Advanced Operating System course and that it will be considered valid up until
the AOS exam of Sept. 2022. 

\begin{flushright}
Group Students' signatures
\end{flushright}


\end{document}